\section{Class architecture}\label{sec:arch-classes}

% ─────────────────────────────────────────────────────────────────────────────
\subsection{Products and cash flows}\label{sec:arch-products}
\code{Product} and the instruments; \code{flows/} coupons and legs; generators.

\pending

Equity options derive from \code{Option}, which holds the option type, the strike (a float64 variable), the
exercise type and the underlying. \code{VanillaOption} is European or American; \code{DigitalOption} is European, with a
\code{payoff\_type} (\code{CashOrNothing} or \code{AssetOrNothing}, enum \code{PayoffType}) and a \code{cash\_amount}.
\code{BlackScholesPricer} prices both, the digitals through \code{digital\_calc}; the factory maps a digital to it.
\code{VanillaMCPricer} prices both too; for the digital payoff \code{digital\_smoothing} sets the sigmoid width (default
1\% of the strike, 0 for the hard indicator). \nb{digital\_options.ipynb}

The barrier option class, \code{BarrierOption} in \code{instruments/barrier.py}, adds the level \code{barrier}, the
type \code{barrier\_type} (an enum \code{BarrierType} with the members \code{DownIn}, \code{DownOut}, \code{UpIn} and
\code{UpOut}, and the flags \code{is\_down} and \code{is\_in}) and the dates \code{monitoring\_dates}, where \code{None}
means continuous monitoring (property \code{continuous}). \code{BlackScholesPricer} prices it through \code{barrier\_calc}
(\code{pricers/black.py}); \code{VanillaMCPricer} through \code{\_bridge\_survival} (continuous) or
\code{\_discrete\_survival} (monitoring dates merged into the grid).

With the autocall at an unreachable level, an \code{AutocallableOption} is a strip of digitals (put barrier 0) or a
put knocked in at expiry (coupons 0).
\nb{digital\_strip.ipynb}, \nb{down\_and\_in\_put.ipynb}

% ─────────────────────────────────────────────────────────────────────────────
\subsection{Pricers}\label{sec:arch-pricers}
\code{Pricer}: \code{calculate\_price}, \code{price(..., autodiff)}, the tape; \code{PricerAssignment}.

A pricer implements \code{calculate\_price(product, market\_env)}: a TensorFlow graph from the market objects (read only
through the \code{MarketEnvironment} accessors, never copied or detached with \code{.numpy()}) to a float64 price.
\code{Pricer.price(product, market\_env, autodiff=False)} runs it, stores the price on \code{product.price} and, with
\code{autodiff=True}, keeps the \code{GradientTape} on \code{pricer.tape}: the greeks are
\code{pricer.tape.gradient(product.price, variables)} for any market variables (or all of them,
\code{market\_env.risk\_factors()}); a second tape around the call gives second orders (gamma). The only variable a
pricer may create is the time to maturity, left on the product (\code{product.time\_to\_maturity}) for theta.
\code{PricerAssignment.create(product)} picks the default pricer by \code{isinstance}: \code{VanillaOption}
$\to$ \code{BlackScholesPricer}, the rates products to their discounting / Bachelier pricers; Monte Carlo pricers are
chosen explicitly. \code{BlackScholesPricer(dividend\_model, use\_implied\_repo)}: European by the Black formula,
American by Bjerksund--Stensland (\code{pricers/black.py}), inputs from the \code{Equity} view, diagnostics on the
product (Market data, The Equity view). \nb{black\_scholes.ipynb}

% ─────────────────────────────────────────────────────────────────────────────
\subsection{Models}\label{sec:arch-models}
\code{StochasticProcess} (\code{size}, \code{initial\_values}, \code{noise\_size}, \code{evolve}) and the
\code{from\_equity} factories used by the Monte Carlo pricers.

\code{StochasticProcess} (\code{models/stochasticprocess.py}) is the common interface: \code{size()} (state
dimension), \code{initial\_values()}, \code{noise\_size()} (normals per path and step, default 1) and
\code{evolve(t\_grid, dw, observation\_times=None)}, which maps pre-drawn normals \code{dw}
($[n_{\text{paths}}, n_{\text{steps}}]$, or with a trailing axis of size \code{noise\_size()}) to paths on the increasing grid
\code{t\_grid} (no 0), returning only the \code{observation\_times} columns when given. The caller draws the normals
(\code{tf.random.stateless\_normal} with a seed), which makes the paths reproducible and the bumps common-random-number
ones. Equity models have a \code{from\_equity(equity, maturity, strike)} factory that builds them from the market
variables, so the simulation is differentiable w.r.t.\ them. They are the GBM (\code{GeometricBrownianMotion},
exact steps), local vol, displaced diffusion, Heston and Heston SLV.
\code{VanillaMCPricer(model, n\_paths, n\_steps, seed, digital\_smoothing)} takes the model \emph{class}, calls
\code{from\_equity} at each price, simulates on a uniform grid ending at the expiry and leaves
\code{standard\_error}, \code{discount\_factor} and \code{time\_to\_maturity} on the product. \nb{monte\_carlo\_gbm.ipynb}

% ─────────────────────────────────────────────────────────────────────────────
\subsection{Calibration}\label{sec:arch-calibration}
Calibration inside the factories (implicit-function greeks) and the rates calibration module.

\pending
